A principis dels anys seixanta del segle xx, els físics Robert Pound, Glen Anderson Rebka i Joseph Snider van verificar al Jefferson Physical Laboratory de Harvard la predicció d'Einstein que la gravetat canvia la freqüència de la llum. Entendre aquest efecte és essencial per a la navegació moderna, com per exemple per al funcionament del GPS. L'experiment consistia a mesurar la variació de freqüència d'uns fotons entre dos punts a diferent altura.

\figura[0.35]{fig1.png}

\begin{apartat}{1,25}
Calculeu la freqüència, la massa (vegeu la nota) i la quantitat de moviment dels fotons al terra del laboratori de l'experiment si tenen una energia de 14,4~keV.
\end{apartat}

\begin{apartat}{1,25}
L'energia mecànica dels fotons és la suma de l'energia dels fotons i de l'energia potencial gravitatòria. A partir del principi de conservació de l'energia mecànica, calculeu la variació (en valor absolut) de l'energia i de la freqüència dels fotons entre dos punts separats verticalment 22,6~m. És a dir, entre el terra del laboratori i un altre punt a la mateixa vertical, 22,6~m més amunt. En quin punt el fotó té una freqüència més gran, quan es troba al terra o quan està 22,6~m per sobre del terra?
\end{apartat}

\nota{Tot i que la massa en repòs d'un fotó és zero, la seva massa efectiva quan és atret per les altres masses és: $m = \dfrac{E_\textit{fotó}}{c^2}$.}

\dades{%
  $1\ \mathrm{eV} = 1,602\times 10^{-19}\ \mathrm{J}$.\\
  $c = 3,00\times 10^{8}\ \mathrm{m\,s^{-1}}$.\\
  $h = 6,63\times 10^{-34}\ \mathrm{J\,s}$.\\
  $g = 9,81\ \mathrm{m/s^2}$.}
